\chapter{Desain Fisikal: Tipe, Constraint, dan Indeks}
\label{modul:desain-fisikal}

\section{Identitas Modul}

\begin{identitasmodul}
\textbf{Sumber materi}    & Pertemuan 6 \\
\textbf{CPMK}             & CPMK-092 \\
\textbf{Minggu pelaksanaan} & Minggu ke-7 \\
\textbf{Durasi tatap muka} & 120 menit \\
\textbf{Estimasi tugas}   & 4--5 jam di luar sesi \\
\textbf{Moda kerja}       & Individual \\
\textbf{Perangkat}        & PostgreSQL 17 dan DBeaver \\
\textbf{Dataset}          & Warehouse NusaMart hasil Modul 4 \\
\textbf{Skrip pendukung}  & \berkas{sql/modul-05/} beserta \berkas{check.sql} \\
\textbf{Luaran}           & \berkas{ddl-fisikal.sql}, \berkas{kamus-tipe-data.md}, dan \berkas{catatan-explain.md} \\
\end{identitasmodul}

\slideref{6}

\section{Capaian dan Indikator Keberhasilan}

Mahasiswa mampu menurunkan model logikal menjadi struktur fisik, memilih tipe
data berdasarkan kebutuhan dan biaya penyimpanan, memasang constraint, serta
mengukur dampak indeks. Bukti analisis harus membandingkan halaman buffer,
bukan hanya waktu eksekusi.

\begin{capaian}
Pada akhir modul ini mahasiswa mampu:
\begin{enumerate}[leftmargin=1.4em]
  \item menata objek gudang data ke dalam schema \perintah{staging},
        \perintah{gudang}, dan \perintah{mart} dengan konvensi penamaan yang
        konsisten;
  \item menghitung lebar baris tabel fakta sampai ke byte --- termasuk header
        tuple, pointer, alignment, dan padding --- lalu memverifikasinya dengan
        \perintah{pg\_column\_size};
  \item memilih tipe data berdasarkan rentang nilai yang benar-benar
        diperlukan, bukan berdasarkan kebiasaan;
  \item memasang constraint lengkap beserta baris \perintah{unknown} sehingga
        tidak ada foreign key fakta yang \perintah{NULL};
  \item membaca \perintah{EXPLAIN (ANALYZE, BUFFERS)} dan membandingkan dampak
        indeks dalam satuan \emph{halaman}, bukan hanya detik.
\end{enumerate}
\end{capaian}

Sampai Modul 4, setiap keputusan diambil di tingkat logika. Mulai modul ini
setiap keputusan memiliki harga yang dapat dihitung: sebuah tipe data yang
dipilih sembarangan pada tabel berisi lima juta baris berbiaya puluhan
megabyte, dan sebuah indeks yang tidak dipakai berbiaya ruang tanpa memberi apa
pun.

\section{Prasyarat dan Persiapan}

\subsection{Prasyarat}

\begin{itemize}
  \item Ketiga tabel fakta dan enam dimensi hasil Modul 4, seluruhnya terisi;
  \item DDL PostgreSQL: \perintah{ALTER TABLE}, \perintah{CREATE INDEX}, serta
        constraint \perintah{NOT NULL}, \perintah{CHECK}, dan
        \perintah{FOREIGN KEY};
  \item pengertian indeks B-tree dan bedanya dengan \istilah{sequential scan};
  \item kemampuan membaca keluaran \perintah{EXPLAIN} secara umum --- bentuk
        pohon rencana dan urutan pembacaannya.
\end{itemize}

\subsection{Pre-lab}

\begin{enumerate}[leftmargin=1.4em]
  \item Bila Modul 4 belum selesai, pulihkan \berkas{snapshot-modul-04.sql}.
  \item Jalankan \perintah{ANALYZE} pada seluruh tabel agar statistik
        perencana mutakhir. Tanpa ini, seluruh pengukuran EXPLAIN Anda akan
        menyesatkan.
\begin{lstlisting}[style=sqlgaya]
ANALYZE;
SELECT relname, n_live_tup, last_analyze
FROM   pg_stat_user_tables ORDER BY n_live_tup DESC;
\end{lstlisting}
  \item Siapkan jawaban berikut pada \berkas{modul-05-fisikal/pre-lab.md}:
  \begin{enumerate}[label=\alph*.,leftmargin=1.4em]
    \item Harga produk NusaMart dicatat dalam rupiah. Apakah rupiah memiliki
          satuan pecahan yang dipakai di kasir? Apa akibatnya bagi pilihan
          tipe data \perintah{NUMERIC(12,2)}?
    \item Berapa nilai transaksi terbesar yang mungkin terjadi di NusaMart?
          Apakah \perintah{INTEGER} memadai untuk menyimpannya?
    \item Sebutkan satu kolom pada tabel fakta Anda yang tipenya menurut Anda
          terlalu besar, beserta alasannya.
  \end{enumerate}
\end{enumerate}

\section{Konsep Dasar}

\subsection{Schema dan Konvensi Penamaan}

Sampai sekarang seluruh objek berada pada \perintah{public}. Untuk gudang data
yang dipakai bersama, pemisahan berdasarkan peran memberi tiga hal sekaligus:
kejelasan, kemudahan pemberian hak akses, dan batas tanggung jawab yang tegas.

\begin{center}
\small
\begin{tabular}{@{}p{0.16\textwidth}>{\raggedright\arraybackslash}p{0.38\textwidth}>{\raggedright\arraybackslash}p{0.38\textwidth}@{}}
\toprule
\textbf{Schema} & \textbf{Isi} & \textbf{Siapa yang membacanya} \\
\midrule
\perintah{staging} & Salinan sumber apa adanya, boleh kotor & Hanya proses ETL \\
\perintah{gudang}  & Dimensi dan fakta yang sudah bersih & Analis dan proses mart \\
\perintah{mart}    & Agregat dan view siap saji & Dashboard dan pengguna akhir \\
\bottomrule
\end{tabular}
\end{center}

Konvensi penamaan yang dipakai sepanjang praktikum: seluruhnya
\perintah{snake\_case} huruf kecil, tanpa singkatan yang tidak baku, dengan
awalan \perintah{dim\_} dan \perintah{fakta\_} yang menyatakan peran tabel.
Nama kolom kunci berakhiran \perintah{\_key} untuk surrogate key dan
\perintah{\_id} untuk natural key --- pembedaan yang sudah Anda pakai sejak
Modul 2.

\subsection{Tipe Data dan Lebar Baris}

Satu baris PostgreSQL tidak hanya berisi datanya. Ia disimpan di dalam
\istilah{page} berukuran 8192 byte, dan setiap baris membawa ongkos tetap
selain isinya.

\begin{konsep}{Anatomi satu baris}
\begin{itemize}
  \item \textbf{Header tuple} --- 23 byte, dibulatkan ke atas menjadi 24 byte
        oleh \perintah{MAXALIGN} 8 byte.
  \item \textbf{Null bitmap} --- $\lceil n/8 \rceil$ byte, \emph{hanya} ada bila
        baris tersebut benar-benar memuat \perintah{NULL}. Baris yang seluruh
        kolomnya terisi tidak membayar apa pun di sini.
  \item \textbf{Data kolom} --- beserta padding yang timbul dari
        \istilah{alignment}.
  \item \textbf{Line pointer} --- 4 byte per baris, disimpan di kepala page,
        bukan di dalam barisnya.
\end{itemize}
Satu page menyediakan $8192 - 24 = 8168$ byte untuk pasangan
pointer dan baris.
\end{konsep}

Padding lahir dari alignment: kolom \perintah{BIGINT} harus dimulai pada offset
kelipatan 8, \perintah{INTEGER} pada kelipatan 4. Bila urutan kolom memaksa
lompatan, byte di antaranya terbuang.

\begin{center}
\small
\begin{tabular}{@{}p{0.24\textwidth}p{0.14\textwidth}p{0.12\textwidth}>{\raggedright\arraybackslash}p{0.42\textwidth}@{}}
\toprule
\textbf{Tipe} & \textbf{Ukuran} & \textbf{Align} & \textbf{Catatan} \\
\midrule
\perintah{BOOLEAN}   & 1 byte  & 1 & \\
\perintah{SMALLINT}  & 2 byte  & 2 & sampai $\pm$32.767 \\
\perintah{INTEGER}   & 4 byte  & 4 & sampai $\pm$2,1 miliar \\
\perintah{BIGINT}    & 8 byte  & 8 & \\
\perintah{DATE}      & 4 byte  & 4 & \\
\perintah{TIMESTAMP} & 8 byte  & 8 & \\
\perintah{NUMERIC}   & berubah & 4 & 2 byte per 4 angka, ditambah 3--8 byte
  ongkos \\
\perintah{VARCHAR(n)} & berubah & 1 & 1 byte ongkos bila isi $\le$126 byte \\
\bottomrule
\end{tabular}
\end{center}

\begin{peringatan}
\perintah{NUMERIC} adalah tipe termahal pada tabel fakta, dan paling sering
dipilih karena kebiasaan. Ia diperlukan ketika presisi desimal wajib dijaga
--- tetapi rupiah tidak memiliki satuan pecahan yang dipakai di kasir.
Menyimpan harga sebagai \perintah{NUMERIC(12,2)} berarti membayar sekitar
sebelas byte untuk menyimpan dua angka di belakang koma yang selalu nol.
\end{peringatan}

Perkiraan di atas cukup untuk merancang, tetapi bukan kebenaran. Ukuran
sebenarnya bergantung pada nilai yang tersimpan, dan PostgreSQL bersedia
memberitahukannya:

\begin{lstlisting}[style=sqlgaya]
SELECT pg_column_size(f.*) AS byte_per_baris
FROM   gudang.fakta_penjualan f LIMIT 5;
\end{lstlisting}

Inilah pola kerja modul ini: hitung dengan tangan, lalu verifikasi dengan
basis data, lalu jelaskan selisihnya.

\subsection{Constraint dan Baris Unknown}

Constraint pada gudang data bukan sekadar warisan kebiasaan OLTP. Ia
menjalankan dua tugas: menolak data yang salah pada saat pemuatan, dan memberi
tahu perencana query fakta yang dapat dipakainya untuk menyusun rencana yang
lebih baik.

Persoalan muncul pada baris fakta yang tidak menemukan pasangan dimensi ---
transaksi tanpa kartu anggota, atau produk yang belum masuk dimensi. Godaan
pertama adalah membiarkan kuncinya \perintah{NULL}.

\begin{konsep}{Mengapa baris unknown, bukan NULL}
Kunci \perintah{NULL} menimbulkan tiga masalah sekaligus:
\begin{enumerate}[leftmargin=1.4em]
  \item \perintah{INNER JOIN} ke dimensi akan membuang baris itu diam-diam,
        sehingga total penjualan menyusut tanpa jejak;
  \item constraint \perintah{NOT NULL} tidak dapat dipasang, sehingga kesalahan
        pemuatan yang sebenarnya tidak lagi terdeteksi;
  \item pengguna akhir tidak memperoleh keterangan apa pun --- ``tidak ada
        baris'' berbeda maknanya dari ``pembeli tanpa kartu anggota''.
\end{enumerate}
Penyelesaiannya adalah satu baris khusus pada setiap dimensi, berkunci
$-1$ dan berlabel \emph{Tidak Diketahui}. Fakta menunjuk baris itu, bukan
\perintah{NULL}.
\end{konsep}

\subsection{Indeks untuk Beban Analitik}

Indeks pada gudang data dipilih dengan pertimbangan yang berbeda dari OLTP.
Query analitik jarang mencari satu baris; ia menyaring rentang lalu
mengagregasi. Karena itu indeks yang berguna adalah yang mempersempit
\emph{rentang} pembacaan.

Urutan kolom pada indeks gabungan menentukan segalanya. Indeks atas
\perintah{(toko\_key, tanggal\_key)} dapat melayani query yang menyaring
\perintah{toko\_key} saja, tetapi \emph{tidak} melayani query yang hanya
menyaring \perintah{tanggal\_key} --- karena kolom pertama indeks tidak
disebut. Aturannya: kolom yang paling sering disaring dengan kesamaan
diletakkan lebih dahulu.

\begin{catatan}
Indeks bukan barang gratis. Setiap indeks menambah ruang dan memperlambat
pemuatan. Pada gudang data yang dimuat sekali sehari, biaya pemuatan itu masih
wajar --- tetapi indeks yang tidak pernah dipakai tetap merupakan pemborosan
murni. Modul ini menuntut bukti bahwa setiap indeks yang Anda pasang benar-benar
dipakai.
\end{catatan}

\subsection{Membaca EXPLAIN ANALYZE BUFFERS}

\perintah{EXPLAIN} menampilkan rencana; \perintah{ANALYZE} menjalankannya dan
melaporkan kenyataan; \perintah{BUFFERS} melaporkan berapa halaman yang
disentuh.

\begin{center}
\small
\begin{tabular}{@{}p{0.28\textwidth}>{\raggedright\arraybackslash}p{0.64\textwidth}@{}}
\toprule
\textbf{Yang dibaca} & \textbf{Artinya} \\
\midrule
\perintah{Seq Scan} lawan \perintah{Index Scan} & Cara tabel dibaca \\
\perintah{rows=} pada \emph{plan} lawan \emph{actual} & Ketepatan perkiraan
  perencana; selisih besar menandakan statistik usang \\
\perintah{shared hit} & Halaman yang ditemukan di cache \\
\perintah{shared read} & Halaman yang benar-benar dibaca dari disk \\
\perintah{Execution Time} & Waktu, yang paling mudah berubah-ubah \\
\bottomrule
\end{tabular}
\end{center}

\begin{peringatan}
Waktu eksekusi adalah ukuran yang paling buruk untuk membandingkan dua
rencana. Eksekusi kedua atas query yang sama hampir selalu lebih cepat karena
halamannya sudah berada di cache --- terlihat sebagai \perintah{shared hit}
yang tinggi dan \perintah{shared read} yang nol. Yang dibandingkan pada modul
ini adalah \textbf{jumlah halaman}, karena angka itu tidak bergantung pada
keadaan cache maupun beban laptop.
\end{peringatan}

\section{Studi Kasus dan Protokol Praktikum}

Skema hasil Modul 4 ditata ulang secara fisik. Setiap keputusan tipe data,
constraint, dan indeks harus disertai alasan dan bukti pengukuran.

Sasaran modul ini adalah \perintah{fakta\_penjualan}, tabel terbesar yang Anda
miliki. Perubahan satu byte pada tabel berisi lima juta baris berarti selisih
lima megabyte --- dan Modul 6 akan bekerja tepat pada dataset sebesar itu.

\begin{catatanasisten}
Pengukuran EXPLAIN pada dataset kecil akan sering memenangkan
\emph{sequential scan}, dan itu bukan kesalahan mahasiswa: pada 100 ribu baris
membaca seluruh tabel memang lebih murah daripada menelusuri indeks. Tegaskan
hal ini sebagai temuan yang sah, dan kaitkan dengan Modul 6 yang memakai
dataset sedang. Yang dinilai adalah kemampuan membaca plan, bukan keberhasilan
memenangkan indeks.
\end{catatanasisten}

\section{Alur Praktikum 120 Menit}

\begin{alurwaktu}
0--15 & Demonstrasi lebar baris dan pembacaan execution plan. \\
15--95 & Kerja terbimbing: schema, tipe, constraint, indeks, dan EXPLAIN. \\
95--110 & Checkpoint DDL fisik dan perbandingan buffer. \\
110--120 & Penjelasan proyeksi ukuran lima tahun. \\
\end{alurwaktu}

\section{Prosedur Praktikum Terbimbing}

\subsection{Bagian A: Menata Schema dan Nama Objek}

\langkah{A.1 Membuat tiga schema dan memindahkan objek}{10 menit}

\begin{lstlisting}[style=sqlgaya]
CREATE SCHEMA IF NOT EXISTS staging;
CREATE SCHEMA IF NOT EXISTS gudang;
CREATE SCHEMA IF NOT EXISTS mart;

-- Objek berpindah tanpa disalin: hanya katalog yang berubah.
ALTER TABLE staging_penjualan SET SCHEMA staging;
ALTER TABLE staging_produk    SET SCHEMA staging;
-- ... ulangi untuk seluruh tabel staging

ALTER TABLE dim_tanggal     SET SCHEMA gudang;
ALTER TABLE dim_produk      SET SCHEMA gudang;
ALTER TABLE fakta_penjualan SET SCHEMA gudang;
-- ... ulangi untuk seluruh dimensi dan fakta
\end{lstlisting}

Perintah \perintah{SET SCHEMA} hanya mengubah katalog, tidak memindahkan satu
byte pun data --- karena itu ia selesai seketika bahkan pada tabel besar.

\begin{catatan}
Sesudah pemindahan, query tanpa nama schema akan gagal. Untuk kenyamanan
selama praktikum, atur \perintah{search\_path}; tetapi pada skrip yang
dikumpulkan, tulis nama lengkap objek. Skrip yang bergantung pada
\perintah{search\_path} tidak reproduktif di lingkungan orang lain.
\end{catatan}

\begin{lstlisting}[style=sqlgaya]
SET search_path TO gudang, staging, mart, public;
\end{lstlisting}

\langkah{A.2 Memberi schema mart isi pertamanya}{5 menit}

\begin{lstlisting}[style=sqlgaya]
CREATE VIEW mart.v_penjualan_bulanan AS
SELECT dt.tahun, dt.tahun_bulan, dp.kategori, dtk.provinsi,
       SUM(f.kuantitas) AS unit_terjual,
       SUM(f.subtotal)  AS nilai_penjualan
FROM   gudang.fakta_penjualan f
JOIN   gudang.dim_tanggal dt  ON dt.tanggal_key = f.tanggal_key
JOIN   gudang.dim_produk  dp  ON dp.produk_key  = f.produk_key
JOIN   gudang.dim_toko    dtk ON dtk.toko_key   = f.toko_key
GROUP  BY 1, 2, 3, 4;
\end{lstlisting}

View ini menjadi titik awal Modul 9. Pada Modul 6 Anda akan membandingkannya
dengan \istilah{aggregate table} dan \istilah{materialized view} untuk
pertanyaan yang sama.

\subsection{Bagian B: Memilih Tipe dan Menghitung Lebar Baris}

\langkah{B.1 Menghitung lebar baris dengan tangan}{15 menit}

Lengkapi Tabel~\ref{tab:m5-lebar-baris} untuk rancangan \emph{sesudah}
perbaikan tipe. Kolom \emph{offset} diisi berurutan, dan padding dicatat
setiap kali sebuah kolom harus melompat ke batas alignment-nya.

\begin{table}[htbp]
\centering
\small
\caption{Kerangka perhitungan lebar baris \perintah{gudang.fakta\_penjualan}}
\label{tab:m5-lebar-baris}
\begin{tabular}{@{}p{0.22\textwidth}p{0.16\textwidth}rrr@{}}
\toprule
\textbf{Kolom} & \textbf{Tipe} & \textbf{Align} & \textbf{Byte} &
  \textbf{Offset} \\
\midrule
\emph{header tuple} & --- & 8 & 23 & 0 \\
\emph{padding}      & --- & --- & 1 & 23 \\
\perintah{transaksi\_id} & \perintah{BIGINT}  & 8 & 8 & 24 \\
\perintah{subtotal}      & \perintah{BIGINT}  & 8 & 8 & 32 \\
\perintah{tanggal\_key}  & \perintah{INTEGER} & 4 & 4 & 40 \\
\perintah{produk\_key}   & \perintah{INTEGER} & 4 & 4 & 44 \\
\perintah{toko\_key}     & \perintah{INTEGER} & 4 & 4 & 48 \\
\perintah{flag\_key}     & \perintah{INTEGER} & 4 & 4 & 52 \\
\perintah{harga\_satuan} & \perintah{INTEGER} & 4 & 4 & 56 \\
\perintah{diskon}        & \perintah{INTEGER} & 4 & 4 & 60 \\
\perintah{kuantitas}     & \perintah{NUMERIC} & 4 & \ldots & 64 \\
\perintah{nomor\_struk}  & \perintah{VARCHAR} & 1 & \ldots & \ldots \\
\emph{padding akhir}     & --- & 8 & \ldots & \ldots \\
\midrule
\multicolumn{3}{@{}l}{\textbf{Lebar baris}} & \ldots & \\
\multicolumn{3}{@{}l}{\textbf{Ditambah line pointer}} & $+\,4$ & \\
\bottomrule
\end{tabular}
\end{table}

Dari lebar baris, hitung tiga angka berikut:

\begin{enumerate}[leftmargin=1.4em]
  \item baris per page $= \lfloor 8168 / (\text{lebar baris} + 4) \rfloor$;
  \item jumlah page $= \lceil \text{jumlah baris} / \text{baris per page}
        \rceil$;
  \item ukuran heap $= \text{jumlah page} \times 8$ KB.
\end{enumerate}

\langkah{B.2 Memverifikasi dengan basis data}{10 menit}

\begin{lstlisting}[style=sqlgaya]
-- Lebar baris sebenarnya, beberapa contoh.
SELECT MIN(pg_column_size(f.*)) AS terkecil,
       ROUND(AVG(pg_column_size(f.*)), 1) AS rata_rata,
       MAX(pg_column_size(f.*)) AS terbesar
FROM   gudang.fakta_penjualan f;

-- Ukuran objek sebenarnya.
SELECT pg_size_pretty(pg_relation_size('gudang.fakta_penjualan'))       AS heap,
       pg_size_pretty(pg_indexes_size('gudang.fakta_penjualan'))        AS indeks,
       pg_size_pretty(pg_total_relation_size('gudang.fakta_penjualan')) AS total;
\end{lstlisting}

Bandingkan hasil hitungan tangan dengan angka ini. Selisih beberapa persen
adalah wajar dan berasal dari \istilah{fillfactor} serta ruang bebas di dalam
page. Selisih puluhan persen berarti ada yang keliru pada perhitungan Anda ---
paling sering pada padding atau pada perkiraan ukuran \perintah{NUMERIC}.

\langkah{B.3 Menerapkan tipe yang benar}{10 menit}

Tulis \berkas{ddl-fisikal.sql} yang membangun ulang tabel fakta dengan tipe dan
urutan kolom yang sudah diperbaiki. Untuk setiap perubahan, catat alasannya
pada \berkas{kamus-tipe-data.md}.

\begin{center}
\small
\begin{tabular}{@{}p{0.20\textwidth}p{0.19\textwidth}p{0.16\textwidth}>{\raggedright\arraybackslash}p{0.37\textwidth}@{}}
\toprule
\textbf{Kolom} & \textbf{Sebelum} & \textbf{Sesudah} & \textbf{Alasan} \\
\midrule
\perintah{subtotal} & \perintah{NUMERIC(14,2)} & \perintah{BIGINT} &
  Rupiah tanpa pecahan; nilai dapat melampaui 2,1 miliar \\
\perintah{harga\_satuan} & \perintah{NUMERIC(12,2)} & \perintah{INTEGER} &
  Rupiah tanpa pecahan; harga satuan jauh di bawah 2,1 miliar \\
\perintah{diskon} & \perintah{NUMERIC(12,2)} & \perintah{INTEGER} & idem \\
\perintah{kuantitas} & \perintah{NUMERIC(10,2)} & \ldots &
  Adakah produk yang dijual dalam pecahan? \\
\bottomrule
\end{tabular}
\end{center}

\begin{peringatan}
Mengubah satuan uang dari rupiah-dengan-sen menjadi rupiah-bulat adalah
keputusan bisnis, bukan sekadar keputusan teknis. Sebelum menerapkannya,
buktikan dengan query bahwa tidak ada satu pun nilai pada
\perintah{staging\_penjualan} yang memiliki angka di belakang koma bukan nol.
Bila ada, catat berapa barisnya dan bahas akibatnya.
\end{peringatan}

\subsection{Bagian C: Memasang Constraint dan Baris Unknown}

\langkah{C.1 Membuat baris unknown}{10 menit}

Jalankan \berkas{sql/modul-05/\allowbreak 02-baris-unknown.sql}, lalu perhatikan
klausa \perintah{OVERRIDING SYSTEM VALUE}: ia diperlukan karena surrogate key
dimensi memakai \perintah{GENERATED ALWAYS AS IDENTITY}, yang secara asali
menolak nilai yang ditentukan sendiri.

\begin{lstlisting}[style=sqlgaya]
INSERT INTO gudang.dim_produk
  (produk_key, produk_id, sku, nama_produk, kategori,
   mulai_berlaku, selesai_berlaku, baris_kini, versi)
OVERRIDING SYSTEM VALUE
VALUES (-1, -1, 'TIDAK-DIKETAHUI', 'Tidak Diketahui', 'Tidak Diketahui',
        DATE '1900-01-01', DATE '9999-12-31', TRUE, 1);
\end{lstlisting}

\langkah{C.2 Mengarahkan kunci NULL ke baris unknown}{10 menit}

\begin{lstlisting}[style=sqlgaya]
UPDATE gudang.fakta_penjualan SET produk_key = -1 WHERE produk_key IS NULL;
UPDATE gudang.fakta_penjualan SET toko_key   = -1 WHERE toko_key   IS NULL;

ALTER TABLE gudang.fakta_penjualan ALTER COLUMN produk_key SET NOT NULL;
ALTER TABLE gudang.fakta_penjualan ALTER COLUMN toko_key   SET NOT NULL;
\end{lstlisting}

\langkah{C.3 Memasang constraint lengkap dan mengujinya}{10 menit}

\begin{lstlisting}[style=sqlgaya]
ALTER TABLE gudang.fakta_penjualan
  ADD CONSTRAINT ck_kuantitas_positif CHECK (kuantitas > 0),
  ADD CONSTRAINT ck_harga_positif     CHECK (harga_satuan > 0),
  ADD CONSTRAINT ck_diskon_wajar      CHECK (diskon >= 0
                                        AND diskon <= kuantitas * harga_satuan);
\end{lstlisting}

Untuk setiap constraint, jalankan satu perintah yang \emph{harus} ditolak, lalu
simpan pesan galatnya. Sebagaimana Modul 3, bukti bahwa aturan ditolak lebih
meyakinkan daripada pernyataan bahwa aturan dipatuhi.

\begin{peringatan}
Bila \perintah{ADD CONSTRAINT} gagal, berarti data Anda sudah melanggar aturan
itu sejak Modul 2 --- ingat temuan kuantitas tidak positif dan diskon negatif
pada profiling Modul 1. Jangan melonggarkan constraint agar perintah berhasil.
Hitung dulu berapa baris yang melanggar, putuskan penanganannya, catat
keputusannya, baru pasang constraint-nya.
\end{peringatan}

\subsection{Bagian D: Mengukur Dampak Indeks}

\langkah{D.1 Mengukur keadaan sebelum indeks}{10 menit}

Pilih satu query yang mewakili beban kerja nyata --- misalnya penjualan satu
provinsi pada satu rentang bulan --- lalu ukur.

\begin{lstlisting}[style=sqlgaya]
EXPLAIN (ANALYZE, BUFFERS, FORMAT TEXT)
SELECT dtk.provinsi, dt.tahun_bulan, SUM(f.subtotal)
FROM   gudang.fakta_penjualan f
JOIN   gudang.dim_toko    dtk ON dtk.toko_key   = f.toko_key
JOIN   gudang.dim_tanggal dt  ON dt.tanggal_key = f.tanggal_key
WHERE  dtk.provinsi = 'Lampung'
  AND  dt.tanggal BETWEEN DATE '2024-01-01' AND DATE '2024-06-30'
GROUP  BY 1, 2;
\end{lstlisting}

Catat pada \berkas{catatan-explain.md}: jenis scan, \perintah{shared hit},
\perintah{shared read}, jumlah baris rencana lawan kenyataan, dan waktu.
Jalankan \textbf{dua kali}, dan catat keduanya --- eksekusi kedua akan
memperlihatkan \perintah{shared read} yang jauh lebih kecil karena halamannya
sudah berada di cache.

\langkah{D.2 Memasang indeks dan mengukur ulang}{10 menit}

\begin{lstlisting}[style=sqlgaya]
CREATE INDEX ix_fakta_penjualan_tanggal
  ON gudang.fakta_penjualan (tanggal_key);

CREATE INDEX ix_fakta_penjualan_toko_tanggal
  ON gudang.fakta_penjualan (toko_key, tanggal_key);

ANALYZE gudang.fakta_penjualan;
\end{lstlisting}

Ulangi pengukuran D.1, lalu lengkapi tabel sebelum--sesudah dalam satuan
halaman \emph{dan} detik. Sertakan juga ukuran indeks yang Anda bayar dengan
\perintah{pg\_indexes\_size}.

\begin{catatan}
Bila perencana tetap memilih \perintah{Seq Scan} sesudah indeks dipasang,
jangan memaksanya. Tuliskan mengapa: pada 100 ribu baris, membaca seluruh
tabel memang lebih murah daripada menelusuri indeks lalu mengambil baris satu
per satu. Ini temuan yang benar, bukan kegagalan --- dan Modul 6 akan
mengulangi pengukuran yang sama pada lima juta baris, tempat kesimpulannya
berubah.

Untuk membuktikan bahwa indeksnya memang dapat dipakai, Anda boleh memaksa
sementara dengan \perintah{SET enable\_seqscan = off}, mencatat plan-nya, lalu
mengembalikannya ke \perintah{on}. Jangan meninggalkan setelan itu menyala.
\end{catatan}

\begin{luaran}
\berkas{modul-05-fisikal/ddl-fisikal.sql}, \berkas{kamus-tipe-data.md} berisi
perhitungan lebar baris beserta verifikasinya, \berkas{catatan-explain.md}
berisi tabel sebelum--sesudah, dan \berkas{bukti-constraint.md}.
\end{luaran}

\begin{checkpoint}
Asisten menjalankan \berkas{sql/modul-05/check.sql} dan memeriksa butir berikut:
\begin{ceklis}
  \item ketiga schema ada, dan tidak ada lagi tabel gudang pada
        \perintah{public};
  \item setiap dimensi memiliki baris unknown berkunci $-1$;
  \item tidak ada satu pun foreign key fakta yang \perintah{NULL};
  \item seluruh kunci dimensi pada fakta sudah \perintah{NOT NULL};
  \item constraint \perintah{CHECK} terpasang pada tabel fakta;
  \item indeks terpasang, dan statistik pemakaiannya tercatat;
  \item \berkas{kamus-tipe-data.md} memuat perhitungan lebar baris beserta
        header, pointer, alignment, dan padding;
  \item \berkas{catatan-explain.md} membandingkan halaman, bukan hanya detik,
        dan menyebutkan pengaruh cache pada eksekusi kedua.
\end{ceklis}
\end{checkpoint}

\section{Latihan Mandiri di Kelas}

\latihan{Memprediksi rencana sebelum menjalankannya}

Usulkan satu indeks tambahan untuk query yang belum Anda ukur. Sebelum
membuatnya, tuliskan prediksi Anda: jenis scan apa yang akan dipakai, dan
kira-kira berapa halaman yang akan dibaca.

Baru setelah prediksi tertulis, buat indeksnya dan jalankan
\perintah{EXPLAIN (ANALYZE, BUFFERS)}. Bandingkan prediksi dengan kenyataan,
lalu jelaskan selisihnya. Prediksi yang meleset dengan penjelasan yang tepat
bernilai lebih tinggi daripada prediksi yang benar tanpa alasan.

\latihan{Menakar biaya urutan kolom}

Susun ulang kolom \perintah{fakta\_penjualan} ke urutan terburuk yang dapat
Anda bayangkan --- selang-seling antara \perintah{BIGINT} dan
\perintah{INTEGER}, dengan kolom \perintah{VARCHAR} di tengah.

\begin{enumerate}[leftmargin=1.4em]
  \item Hitung lebar barisnya dengan tangan, lalu bandingkan dengan urutan
        yang Anda pakai.
  \item Berapa byte per baris yang terbuang menjadi padding?
  \item Pada dataset sedang berisi lima juta baris, berapa megabyte selisihnya?
  \item Buat tabel uji dengan urutan buruk itu, isi seribu baris, lalu
        buktikan selisihnya dengan \perintah{pg\_column\_size}.
\end{enumerate}

\section{Tugas Individu}

\subsection{Pekerjaan Wajib}
Hitung proyeksi ukuran tabel untuk lima tahun, termasuk heap, indeks, dan
asumsi pertumbuhan data.

Kumpulkan \berkas{proyeksi-ukuran.md} yang memuat:

\begin{enumerate}[leftmargin=1.4em]
  \item asumsi pertumbuhan yang dinyatakan secara eksplisit --- jumlah
        transaksi harian, jumlah produk per struk, dan laju pertambahan toko;
  \item proyeksi ukuran heap ketiga tabel fakta per tahun sampai tahun kelima;
  \item proyeksi ukuran indeks, dengan asumsi yang Anda nyatakan sendiri
        mengenai rasio indeks terhadap heap;
  \item proyeksi ukuran \perintah{fakta\_persediaan\_harian}, yang tumbuh
        menurut jumlah produk dikalikan jumlah toko dikalikan jumlah hari ---
        perhatikan bahwa ia tidak bergantung pada banyaknya transaksi;
  \item satu paragraf yang menyebutkan pada tahun ke berapa ukurannya menjadi
        persoalan, dan tindakan apa yang perlu disiapkan sebelum itu.
\end{enumerate}

\begin{catatan}
Butir keempat biasanya mengejutkan. Periodic snapshot tumbuh menurut perkalian
tiga besaran, dan pertumbuhannya tidak melambat ketika penjualan sepi. Pada
kebanyakan perhitungan mahasiswa, tabel inilah yang lebih dahulu menjadi
persoalan --- bukan tabel penjualan. Tindakan yang perlu disiapkan adalah
partisi dan agregat, dan keduanya menjadi bahan Modul 6.
\end{catatan}

\subsection{Pertanyaan Analisis}

\begin{enumerate}[leftmargin=1.4em]
  \item Anda memilih \perintah{INTEGER} untuk \perintah{harga\_satuan}. Lima
        tahun kemudian NusaMart membuka lini produk elektronik seharga puluhan
        juta rupiah. Apakah pilihan Anda masih memadai? Tunjukkan
        perhitungannya, dan sebutkan apa yang harus dilakukan bila tidak.
  \item Seorang rekan mengusulkan menghapus seluruh foreign key pada tabel
        fakta agar pemuatan lebih cepat. Sebutkan apa yang diperoleh dan apa
        yang hilang, lalu berikan pendirian Anda beserta alasannya.
  \item Indeks atas \perintah{(toko\_key, tanggal\_key)} sudah ada. Seorang
        rekan hendak menambah indeks atas \perintah{(toko\_key)} saja. Jelaskan
        mengapa indeks kedua itu hampir selalu mubazir.
\end{enumerate}

\section{Format Pengumpulan dan Checklist}

Kumpulkan DDL, kamus tipe data, catatan EXPLAIN, dan proyeksi ukuran. Pastikan
perhitungan mencakup header, pointer, alignment, dan padding.

Seluruh berkas diletakkan pada \berkas{modul-05-fisikal/}.

\begin{ceklis}
  \item \berkas{pre-lab.md}
  \item \berkas{ddl-fisikal.sql}
  \item \berkas{kamus-tipe-data.md} --- perhitungan tangan dan verifikasinya
  \item \berkas{catatan-explain.md} --- tabel sebelum--sesudah dalam halaman
        dan detik
  \item \berkas{bukti-constraint.md}
  \item \berkas{proyeksi-ukuran.md}
  \item \berkas{jawaban-analisis.md}
\end{ceklis}

\section{Troubleshooting}

\begin{table}[htbp]
\centering
\small
\caption{Masalah yang lazim muncul pada Modul 5 dan penanganannya}
\label{tab:m5-troubleshooting}
\begin{tabular}{@{}p{0.30\textwidth}>{\raggedright\arraybackslash}p{0.62\textwidth}@{}}
\toprule
\textbf{Gejala} & \textbf{Penanganan} \\
\midrule
\perintah{relation "fakta\_penjualan" does not exist} sesudah pemindahan &
  Objek kini berada pada schema \perintah{gudang}. Tulis nama lengkapnya, atau
  atur \perintah{search\_path} untuk sesi ini. \\
\perintah{cannot insert a non-DEFAULT value into column "produk\_key"} &
  Kolom memakai \perintah{GENERATED ALWAYS AS IDENTITY}. Tambahkan
  \perintah{OVERRIDING SYSTEM VALUE} pada \perintah{INSERT} baris unknown. \\
\perintah{check constraint ... is violated by some row} &
  Data sudah melanggar aturan sebelum constraint dipasang. Hitung barisnya,
  putuskan penanganannya, catat keputusannya --- jangan melonggarkan
  constraint. \\
\perintah{column "produk\_key" contains null values} &
  Masih ada kunci \perintah{NULL}. Arahkan ke baris unknown lebih dahulu,
  baru pasang \perintah{NOT NULL}. \\
Indeks dibuat tetapi tetap \perintah{Seq Scan} &
  Pada dataset kecil ini lazim dan benar. Catat sebagai temuan; bila perlu
  buktikan indeksnya dapat dipakai dengan \perintah{SET enable\_seqscan = off}
  sementara. \\
Perkiraan baris pada plan jauh meleset &
  Statistik usang. Jalankan \perintah{ANALYZE} pada tabel yang bersangkutan,
  lalu ukur ulang. \\
\perintah{shared read} nol pada pengukuran pertama &
  Tabel sudah berada di cache dari pekerjaan sebelumnya. Sebutkan hal ini pada
  catatan; jangan menyimpulkan bahwa query tidak menyentuh disk. \\
Hitungan lebar baris meleset jauh dari
  \perintah{pg\_column\_size} &
  Paling sering karena padding terlewat, atau karena \perintah{NUMERIC}
  diperkirakan sebagai lebar tetap padahal ukurannya bergantung nilai. \\
\bottomrule
\end{tabular}
\end{table}

\section{Ringkasan}

Modul ini mengubah rancangan menjadi objek yang harganya dapat dihitung. Tiga
gagasan perlu dibawa ke Modul 6.

\emph{Pertama}, tipe data adalah keputusan berbiaya. \perintah{NUMERIC} yang
dipilih karena kebiasaan berbiaya beberapa byte per baris, dan pada tabel
berisi jutaan baris beberapa byte itu menjadi puluhan megabyte --- ditambah
setiap halaman tambahan yang harus dibaca oleh setiap query.

\emph{Kedua}, baris \perintah{unknown} bukan kerapian melainkan syarat
kebenaran. Kunci \perintah{NULL} membuat \perintah{INNER JOIN} membuang baris
diam-diam, dan angka yang hilang tanpa jejak jauh lebih berbahaya daripada
angka yang salah secara mencolok.

\emph{Ketiga}, klaim tentang kinerja harus dibuktikan dalam satuan yang tidak
berubah-ubah. Detik bergantung pada cache dan beban laptop; halaman tidak.
Setiap tindakan optimisasi harus dapat menjawab pertanyaan ``berapa halaman
yang dihemat, dan berapa ruang yang dibayar''.

Pada dataset kecil, sebagian besar pengukuran Anda hari ini akan memenangkan
sequential scan. Modul 6 mengulangi pengukuran yang sama pada dataset sedang
berisi lima juta baris --- di sanalah partisi, aggregate table, dan
materialized view mulai memperlihatkan apa yang sebenarnya mereka beli.
